\documentclass[10pt,a4paper]{amsart}
\usepackage[margin=19mm]{geometry}
\usepackage{amsmath,amssymb,mathtools}
\usepackage[scr=rsfs,cal=euler]{mathalfa}
\usepackage{xfrac}
\usepackage[unicode,hidelinks,bookmarks=false]{hyperref}
\newcommand{\ie}{i.e.\ }
\newcommand{\cf}{cf.\ }
\newcommand{\resp}{respectively\ }
\newcommand{\Subsection}{\subsection}
\providecommand{\nobreakdash}{\nobreak\hbox{-}}
\setlength{\emergencystretch}{2em}
\multlinegap0pt
\usepackage{bm}
\usepackage{bbm}
\usepackage{dsfont}
\usepackage{xspace}
\usepackage{cancel}
\usepackage{mathtools}
\usepackage{amsthm}
\makeatletter
\@ifundefined{thm}{\newtheorem{thm}{Theorem}}{}
\makeatother
\DeclareMathOperator{\Sym}{Sym}
\newcommand{\qsres}{quasi-semiregular elements}
\newcommand{\qsre}{quasi-semiregular element}
\newcommand{\qsr}{quasi-semiregular}
\title[Finite permutation groups with quasi-semiregular elements]{Finite permutation groups with quasi-semiregular elements}
\author{Michael Giudici, Luke Morgan, Cheryl E. Praeger}
\date{Source: arXiv:2507.12866v1 (CC BY 4.0)}
\begin{document}
\maketitle

\noindent\emph{Source and licence.} Finite permutation groups with quasi-semiregular elements. Version arXiv:2507.12866v1 (CC BY 4.0), distributed under the Creative Commons Attribution 4.0 International licence, as recorded in the arXiv licence metadata for this exact version and shown on the abstract page https://arxiv.org/abs/2507.12866v1. Full rights notice: licenses/arxiv-2507.12866-CC-BY-4.0.txt.\par\smallskip

\noindent\textit{Editorial scope.} Counted unit: the Thm whose source label is \texttt{t:pa} in the article (source0.tex lines 521--523, source label \texttt{t:pa}), delivered together with the article's own complete original proof of it (source0.tex lines 524--549, one \texttt{proof} environment of 26 lines). No mathematics is written, repaired or re-derived: statement and proof are the article's own text line for line. Required context retained uncounted: none. Every definition, display and label of those lines is retained unchanged. Omitted from this package and not redistributed: 1--520, 550--1394. Nothing inside a retained excerpt is omitted or altered: every retained body is delivered exactly as the article's lines are written, so no omission receipt of any kind accompanies this package. Citations: the retained excerpts carry no citation command at all, so no bibliography entry of the article is transcribed and no reference is redistributed. Reference adaptation: none was needed and none was made; every reference inside the delivered unit resolves inside it, so no unresolved reference or citation is produced. Printed slips kept literal: no printed word, symbol, hypothesis or number of the counted statement or of its proof was altered. Layout and class: the article's own class is \texttt{article} with packages including amsmath, amssymb, amsthm, a4wide, color, enumerate, cleveref, mathrsfs, url; the wrapper is the lane's pinned \texttt{amsart} editorial wrapper at 10pt on A4, which restates the theorem-like environments the retained text needs and adds the article's own macro definitions that the retained text needs (\texttt{\textbackslash Sym}, \texttt{\textbackslash qsr}, \texttt{\textbackslash qsre}, \texttt{\textbackslash qsres}), with extra wrapper packages (bm, bbm, dsfont, xspace, cancel, mathtools). The delivered numbering is the unit's own. Assets: no retained span contains a figure environment, a graphics-inclusion command, a PDF-page inclusion, a table or a TikZ picture; no image file is copied into this package, the package declares no asset dependency and no image file is redistributed with it. Licence: version arXiv:2507.12866v1 of this article is distributed under the Creative Commons Attribution 4.0 International licence, as recorded in the arXiv licence metadata for this exact version and shown on the abstract page https://arxiv.org/abs/2507.12866v1. A full rights notice is delivered at licenses/arxiv-2507.12866-CC-BY-4.0.txt. Characters: every retained excerpt is pure ASCII, so the delivered engine, XeLaTeX with its default Latin Modern text font, reports no missing character.
\begin{thm}\label{t:pa}
Let $G= \Sym(k)\wr\Sym(\ell)$ in its natural action on $[k^\ell]$, where $k\geqslant3$ and $\ell\geqslant2$. Then $G$ contains {\qsres}. Moreover, if $g\in G$ is a {\qsre} of prime order, then  $g=(h_1,\ldots,h_\ell)\in \Sym(k)^\ell$ and   $h_i\in \Sym(k)$ is  {\qsr} for all $i$ such that $1\leqslant i\leqslant \ell$.
\end{thm}

\begin{proof}
    Set $H=\Sym(k)$. We view elements of the set   $[k^\ell]$   as tuples $(\alpha_1,\ldots,\alpha_\ell)$ with $\alpha_i\in [k]$ for each $i$. We write elements of $G$ as products $h\sigma $ where $h\in H^\ell$ and $\sigma \in \Sym(\ell)$.  The elements of $G$  act on $[k^\ell]$ via the following rule. For $x=(\alpha_1,\ldots,\alpha_\ell)\in [k^\ell]$ and $h=(h_1,\ldots,h_\ell)\in H^\ell$ we have $x^h :=  (\alpha_1^{h_1},\ldots,\alpha_\ell^{h_\ell})$ and $x^\sigma  = (\alpha_{1\sigma^{-1}},\ldots,\alpha_{\ell\sigma^{-1}})$. Furthermore, $\sigma \in \Sym(\ell)$ acts by conjugation on $ H^\ell$ via $h^\sigma  := (h_{1\sigma^{-1}},\ldots,h_{\ell\sigma^{-1}})$.

    For each $i=1,\dots, \ell$, let  $h_i \in H$ be quasi-semiregular with unique fixed point $\alpha_i \in [k]$ (note that such an $h_i$ exists since $H=\Sym(k)$ and $k\geqslant3$). Define $h=(h_1,\ldots,h_\ell)$. Then $h$ fixes $x:=(\alpha_1,\ldots,\alpha_\ell)\in [k^\ell]$.  Let $y = (\beta_1,\ldots,\beta_\ell) \in  [k^\ell]$  such that $y \neq x$. Then there exists $i$ such that  $\beta_i\neq \alpha_i $, and hence $\beta_i^{h_i} \neq \beta_i$ since $\alpha_i$ is the unique point of $[k]$ fixed by $h_i$. Thus $y^h \neq y$ and so  $h$ is a {{\qsre}} in $G$.

    Conversely suppose that $g\in G$ is a {\qsre} of prime order $p$. After conjugating by some element of $G$, we may assume that     $g$ fixes the point $x=(\alpha,\ldots,\alpha)$ for some $\alpha \in [k]$. Write $g=h\sigma$ with $h=(h_1,\ldots,h_\ell) $ and note that $\alpha^{h_i} = \alpha$ for $1\leqslant i \leqslant\ell$. Since $g^p=1$ we have $\sigma^p=1$. In the action of $\sigma$ on $\ell$ points, suppose that $\langle \sigma\rangle$ has $a$ orbits of size $p$ and $b$ orbits of size $1$ (so $ap+b=\ell$). 
    Suppose that $a\geqslant 1$. After relabelling, we may assume  that $\sigma$ contains the $p$-cycle $(1,\ldots,p)$ (on $\ell$ points). 
Now we have:  
\[
1=g^p=h\sigma h \sigma \ldots h\sigma =h h^{\sigma^{-1}} \ldots h^{\sigma^{-(p-1)}}   \sigma^p = h h^{\sigma^{-1}} \ldots h^{\sigma^{-(p-1)}} 
\]
and we note that, for all $i$ such that $1\leqslant i\leqslant p-1$
$$
h^{\sigma^{-i}}=(h_{1 \sigma^i},\ldots,h_{\ell\sigma^{i}})=(h_{1+i},\ldots,h_p,h_1,\ldots,h_i,h_{(p+1)\sigma^i},\ldots,h_{\ell \sigma^i}).
$$
Thus the first entry of the $\ell$-tuple $h h^{\sigma^{-1}} \cdots h^{\sigma^{-(p-1)}}$ is $h_1h_2\cdots h_p$. Since $h h^{\sigma^{-1}} \cdots h^{\sigma^{-(p-1)}} = g^p = 1$, it follows that $h_1h_2\cdots h_p=1$.
Pick  $\beta \neq \alpha$,  set 
    \[y= (\beta,\beta^{h_1},\beta^{h_2},\ldots,\beta^{h_1h_2\cdots h_{p-1}},\alpha,\ldots,\alpha)\]
    and note that  
\begin{align*}
       y^{h\sigma} &= (\beta^{h_1},\beta^{h_1h_2},\ldots,\beta^{h_1 h_2 \cdots h_p},\alpha^{h_{p+1}},\ldots,\alpha^{h_\ell})^\sigma \\
       & = (\beta^{h_1\cdots h_p},\beta^{h_1},\ldots,\beta^{h_1 \cdots h_{p-1}},\alpha,\ldots,\alpha).
\end{align*}
Since $\beta^{h_1\cdots h_p}=\beta$, we have $y^g = y$. Now $\beta \neq \alpha$ so $y\neq x$, and this is a contradiction to $g $ being {\qsr}. It follows that $a=0$ so that $\sigma =1$ and $h=(h_1,\ldots,h_\ell) \in \Sym(k)^\ell$. 
Finally, since $h$ is {\qsr}, for any $\beta \neq \alpha$ we have $(\beta,\alpha,\ldots,\alpha)\neq (\beta,\alpha,\ldots,\alpha)^h = (\beta^{h_1},\alpha,\ldots,\alpha) $ so that $\beta^{h_1}\neq  \beta$. Thus, $h_1$ is {\qsr}. Similarly $h_i$ is {\qsr} for each $i$.
\end{proof}

\end{document}
